\ifdefined\MyWebBook\else
\documentclass[../textbook.tex]{subfiles}
\begin{document}
\fi
\bookchapter{推理计算优化}{ch:inference-algorithms}

推理计算优化围绕执行时延、显存占用与服务吞吐展开。增量解码通过键值缓存消除历史前缀的重复前向；算子融合与片上分块计算减少显存与计算核心之间的数据搬运；推测解码利用计算不对称性在单次大模型前向中并行验证多个候选词元。不同优化路径对系统等价性有明确约束：算子分块在浮点舍入容差内保证严格数学等价，精确推测解码保持原始自回归采样分布，而低精度量化与窗口截断则引入受控的逼近误差。

\bookxref{ch:incremental-decoding-cache}讨论了键值状态的历史复用。在线归一化与分块注意力抑制了显存往返开销，图融合与特化编译消除了运行时调度气泡，精确推测解码则通过接受与残差重采样保障分布一致性。低精度量化误差分析在\bookxref{ch:quantization}展开，请求级动态调度在\bookxref{ch:batching-scheduling}展开。

\section{推理执行成本}

本章主要符号见表\ref{tab:inference-algorithms-symbols}。
\begin{table}[htbp]
\centering
\caption{推理计算优化主要符号}\label{tab:inference-algorithms-symbols}
\begin{tabular}{ll}
\toprule
符号 & 含义\\
\midrule
$T_q,T_k,d_h,d_v$ & 查询数、键数、键头维数和值头维数。\\
$\mat{Q},\mat{K},\mat{V}$ & 单头矩阵，形状分别为 $T_q\times d_h,T_k\times d_h,T_k\times d_v$。\\
$z_j,v_j$ & 某查询对第 $j$ 个键的分数及对应值向量。\\
$m,\ell,\vect{u}$ & 已处理键集合的最大分数、缩放指数和、缩放值加权和。\\
$F,Q_b,R,b$ & 运算量、搬运字节数、有效计算速率和有效带宽。\\
$p,q$ & 给定前缀下的目标采样分布与草稿采样分布。\\
$k,\alpha$ & 单轮草稿候选数及特定假设下的条件接受概率。\\
\bottomrule
\end{tabular}
\end{table}

对特定的执行内核，计算时间下界为 $\frac{F}{R}$，数据传输时间下界为 $\frac{Q_b}{b}$。二者在硬件流水线中可部分重叠，因此 $\max(\frac{F}{R},\frac{Q_b}{b})$ 构成了理论时延的下限基准。端到端实际延迟还包含主机侧下发调度、设备内核启动（Kernel Launch）、显存同步及张量布局重排。小规模张量算子常受限于启动与同步延迟，长序列大矩阵运算则由算力峰值或全局显存带宽主导。

若某项优化仅作用于占基准总耗时比例为 $f$ 的算子，且该部分获得 $s$ 倍加速，其余部分耗时保持不变，则系统端到端加速比由阿姆达尔定律决定：
\begin{equation}
 S=\frac{1}{(1-f)+\frac{f}{s}}.
 \label{eq:kernel-amdahl}
\end{equation}
例如当优化算子占总执行时间的比例 $f=0.4$ 时，即使算子本身取得 $s=4$ 倍加速，系统总加速比上限仅为 $\frac{1}{0.6+0.1}\approx 1.43$。若优化实现引入了额外的跨张量布局转换或显卡流同步，端到端收益将进一步折损。因此算子加速必须置于完整调用链与计算图中综合度量。

\section{分块注意力计算}

\subsection{分块归一化合并}
单查询注意力输出为
\begin{equation}
 \vect{o}=\frac{\sum_j \conste{}^{z_j}v_j}{\sum_j \conste{}^{z_j}},\qquad
 z_j=\frac{\vect{q}^{\mathrm{T}}k_j}{\sqrt{d_h}}+M_j.
 \label{eq:kernel-attention}
\end{equation}
$M_j$ 表示因果掩码或相对位置偏置。将键值序列划分为多个独立块后，各分块局部 Softmax 的归一化分母互不相同，直接相加或等权平均局部输出会导致加权基准错位；分块计算的关键在于维护各分块的局部统计量，并在归约合并时恢复其在全局指数和中的准确质量权重。

对非空有效键集合 $A$，维护局部三元统计量：
\begin{equation}
 m_A=\max_{j\in A}z_j,\quad
 \ell_A=\sum_{j\in A}\conste{}^{z_j-m_A},\quad
 \vect{u}_A=\sum_{j\in A}\conste{}^{z_j-m_A}v_j.
 \label{eq:kernel-summary}
\end{equation}
其中 $m_A,\ell_A$ 为标量，$\vect{u}_A\in\Real^{d_v}$。在线归一化算法通过随分块遍历动态更新最大值与缩放和，消除了传统实现必须先遍历全序列求全局最大值再计算全局分母的二次全局读写访存\autocite{milakov2018online}。

\subsection{分块统计量的合并}
对不相交有效集合 $A,B$，设全局局部最大值为 $m=\max(m_A,m_B)$，统计量合并递推公式为
\begin{align}
 \ell&=\conste{}^{m_A-m}\ell_A+\conste{}^{m_B-m}\ell_B,\\
 \vect{u}&=\conste{}^{m_A-m}\vect{u}_A+\conste{}^{m_B-m}\vect{u}_B.
 \label{eq:kernel-merge}
\end{align}
由恒等式 $\conste{}^{m_A-m}\conste{}^{z_j-m_A}=\conste{}^{z_j-m}$，合并后的 $\ell$ 严格等于并集 $A\cup B$ 的全局缩放指数和，$\vect{u}$ 同理。合并结果精确满足式\eqref{eq:kernel-summary}定义的数学形式，最终注意力输出为 $\vect{o}=\frac{\vect{u}}{\ell}$。\keyconcept{分块在线归一化通过三元状态 $(m,\ell,\vect{u})$ 维护局部的最大值、指数和与值加权和，在单趟片上遍历中恢复全局精确 Softmax，消除了完整注意力矩阵的高带宽显存往返。}

只要在归约过程中严格维护上述三元统计量，即可按任意分块顺序流式遍历键值集合。在精确算术下，合并操作满足结合律，输出结果与分块拓扑完全无关；实际硬件中不同的分块归约顺序仅引入微小的浮点舍入累积差。

若将全掩码空块的最大值形式化置为 $m=-\infty$，浮点减法运算 $(-\infty)-(-\infty)$ 会直接产生非数（NaN）。工程实现必须通过显式的有效性标志（valid flag）解耦空块处理：完全被掩码屏蔽的块\emph{直接跳过统计量更新}；首个包含合法键的块直接初始化三元组。若整行查询均无有效键，全局分母退化为零，系统应依接口规范显式填零或上报异常，杜绝未定义浮点异常向后级网络传递。

\begin{warning}[全掩码分块的浮点 NaN 陷阱]
若将完全被掩码屏蔽的空块最大值形式化置为 $m=-\infty$，两块合并时计算指数差值将触发 $(-\infty)-(-\infty)$ 的未定义浮点运算，直接在片上生成非数（NaN）。工程实现必须通过显式的有效性标志（valid flag）解耦空块判定：全掩码块直接跳过统计量更新；首个合法块直接初始化三元状态；若整行全被屏蔽，显式填零或上报异常，杜绝浮点异常污染后级计算。
\end{warning}

\begin{example}
设分数值为 $(\log2,\log1,\log6)$，对应标量值向量为 $(1,4,5)$。前两个键构成块 $A$，第三个键构成块 $B$。分别计算两块局部统计量：
\begin{align*}
 m_A&=\log2,&\ell_A&=\frac{3}{2},&u_A&=3,\\
 m_B&=\log6,&\ell_B&=1,&u_B&=5.
\end{align*}
两块联合最大值为 $m=\max(\log2,\log6)=\log6$，块 $A$ 相对最大值的缩放权重为 $\conste{\log2-\log6}=\frac{1}{3}$，代入合并公式：
\begin{equation*}
 \ell=\left(\frac{1}{3}\right)\left(\frac{3}{2}\right)+1=\frac{3}{2},\qquad u=\left(\frac{1}{3}\right)3+5=6,\qquad o=\frac{u}{\ell}=4.
\end{equation*}
若直接计算全局 Softmax，输出为 $\frac{2\times 1+1\times 4+6\times 5}{2+1+6}=4$。块 $A$ 与块 $B$ 的局部归一化输出分别为 $2$ 和 $5$；直接相加得 $7$，等权平均得 $3.5$，均偏离真实解。在线合并通过乘法重缩放因子 $\frac{1}{3}$ 精确恢复了各块在全局指数和中的比重。
\end{example}

\begin{algorithm}[单查询的分块精确注意力]
\AlgoInput{查询向量、按块分批加载的键值张量、因果与注意力掩码。}
\AlgoOutput{精确注意力加权值向量 $\vect{o}$。}
\begin{myalgorithmic}[1]
\State 令 $\mathrm{valid}\gets\mathrm{false}$
\For{each 有限输入中的键值块}
  \State 计算有效位置及其分数
  \If{该块包含至少一个有效位置}
    \State 计算局部块状态 $(m_b,\ell_b,\vect{u}_b)$
    \If{$\mathrm{valid}=\mathrm{false}$}
      \State 令 $(m,\ell,\vect{u})\gets(m_b,\ell_b,\vect{u}_b)$ 且 $\mathrm{valid}\gets\mathrm{true}$
    \Else
      \State 利用式\eqref{eq:kernel-merge}合并块状态
    \EndIf
  \EndIf
\EndFor
\If{$\mathrm{valid}=\mathrm{true}$}
  \State \Return $\frac{\vect{u}}{\ell}$
\Else
  \State 执行输入契约规定的全掩码行行为
\EndIf
\end{myalgorithmic}
\end{algorithm}

\subsection{FlashAttention 的存储组织}

标准注意力的显式实现通常先将完整的 $T_q\times T_k$ 注意力分数矩阵写回高带宽显存（HBM），再次读取后执行行 Softmax，最后将完整的注意力概率矩阵写入显存并读出与 $\mat{V}$ 相乘。当 $T_q=T_k=8192$ 时，单头浮点分数矩阵包含 $\num{67108864}$ 个元素；以半精度（FP16/BF16）存储即占用 \qty{128}{\mebi\byte}，32个注意力头将导致中间显存占用激增至 \qty{4}{\gibi\byte}。

FlashAttention 通过切分查询、键与值矩阵，将分块运算与在线归一化状态完全驻留在片上共享内存（SRAM）中，消除了中间概率矩阵在显存与芯片核心间的频繁往返\autocite{dao2022flashattention}。FlashAttention 属于 IO 感知型（IO-Aware）算法优化，其稠密注意力的理论计算复杂度仍严格为 $\mathcal{O}(T_q T_k d_h)$。因果掩码仅在计算调度层跳过完全不可见的下三角分块，并未改变有效注意力覆盖范围内的计算等价性。

\begin{figure}[htbp]
\centering
\begin{tikzpicture}[node distance=7mm,box/.style={draw,rounded corners,align=center,text width=33mm,minimum height=13mm},>=stealth]
\node[box] (hbm) {显存\\$\mat{Q},\mat{K},\mat{V}$ 分块};
\node[box,right=of hbm] (tile) {片上分数块\\点积、缩放、掩码};
\node[box,right=of tile] (state) {在线状态\\$m,\ell,\vect{u}$};
\node[box,below=of state] (out) {写回最终输出\\$\vect{o}=\frac{\vect{u}}{\ell}$};
\draw[->] (hbm)--(tile);\draw[->] (tile)--(state);\draw[->] (state)--(out);
\draw[->] (state.north) to[out=110,in=70,looseness=1.3] node[above,font=\small] {下一键块继续合并} (tile.north);
\end{tikzpicture}
\caption[分块注意力数据搬运流]{分块注意力数据搬运流。完整分数矩阵不必写回高带宽显存，查询仍通过分块在线归一化遍历全部合法键。}\label{fig:kernel-tiling}
\end{figure}

\subsection{块大小、并行度与反向重算}
较大的块可以复用更多数据，却占用更多寄存器和片上存储，降低同时驻留的执行单元数；过小的块则增加读取和循环开销。因此块大小由头维、精度、设备资源和并行映射共同决定，最大块未必是最优块。

表\ref{tab:flashattention-evolution}梳理了 FlashAttention 算法与硬件协同设计的演进历程。从算子分块、循环排布到深层次的异步管线与硬件特化，注意力内核持续逼近现代加速器的理论计算极限。

\begin{table}[htbp]
\centering
\caption{FlashAttention 算法演进对比}\label{tab:flashattention-evolution}
\begin{tabularx}{\linewidth}{@{}lXXX@{}}
\toprule
特性维度 & FlashAttention-1 & FlashAttention-2 & FlashAttention-3 \\
\midrule
代表硬件与年份 & Ampere (A100), 2022 & Ampere/Ada (A100/H100), 2023 & Hopper (H100/H800), 2024 \\
外层/内层循环 & 外层 $K,V$，内层 $Q$（需原子累加写回） & 外层 $Q$，内层 $K,V$（消除跨块输出累加同步） & 外层 $Q$，内层 $K,V$（两阶段异步乒乓管线） \\
Warp 并行划分 & 块内沿列划分（需共享内存中间交换） & 块内沿行划分（消除 Warp 间共享内存通信） & Warp 特化分工（生产者加载与消费者计算解耦） \\
硬件异步机制 & 基础异步拷贝机制 & 优化通用张量核心与寄存器排布 & 深度利用 TMA 硬件单元与异步事务屏障 \\
数值精度支持 & FP16 / BF16 & FP16 / BF16 & FP16 / BF16，原生支持 FP8 张量计算 \\
典型理论算力利用率 & $25\%\sim 40\%$ & $50\%\sim 70\%$ & $75\%\sim 85\%$ \\
\bottomrule
\end{tabularx}
\end{table}

训练时无需保存全部注意力概率，也可以在反向中由保存的行归一化信息和重算分数恢复局部概率。令 $\mat{P}=\operatorname{softmax}(\mat{S})$、$\mat{O}=\mat{P}\mat{V}$，上游梯度为 $\mat{G}_O$，则
\begin{align}
 \mat{G}_V&=\mat{P}^{\mathrm{T}}\mat{G}_O,&\mat{G}_P&=\mat{G}_O\mat{V}^{\mathrm{T}},\\
 (\mat{G}_S)_{ij}&=P_{ij}\left[(\mat{G}_P)_{ij}-\sum_kP_{ik}(\mat{G}_P)_{ik}\right],\\
 \mat{G}_Q&=\frac{\mat{G}_S\mat{K}}{\sqrt{d_h}},&\mat{G}_K&=\frac{\mat{G}_S^{\mathrm{T}}\mat{Q}}{\sqrt{d_h}}
\end{align}
这些关系解释为什么反向可逐块重新计算所需概率。实际实现还须处理掩码、Dropout和浮点累计；重算时使用的随机掩码必须与正向一致。

\subsection{注意力后端的兼容条件}
缩放点积注意力（SDPA）的高层接口封装了统一的计算语义，但底层具体内核后端的派发取决于硬件架构、输入张量形状、数据类型及内存布局。切换至优化内核时，因果掩码对齐、多查询与分组查询头映射、缩放因子及内存步幅（Stride）连续性必须与规范实现严格一致。

在增量解码阶段，单步查询长度为 $1$，而键值缓存覆盖长度为 $T_k$ 的历史前缀。若将 $1\times T_k$ 因果掩码按局部相对索引从零对齐，会导致仅首个历史键可见而后续历史键全部被错误屏蔽。因果掩码必须依据当前查询在序列中的全局逻辑位置 $T_k-1$ 展开可见性判定。后端集成测试须逐项核验逻辑时间戳与实际注意力可见集合的映射关系。

\section{执行图优化}

执行图优化通过消除中间临时张量的显存往返、减少内核启动频率并复用静态执行拓扑来降低端到端时延。算子融合重构指令流水线，图编译生成硬件特化机器码，图回放消除 CPU 侧的主机下发开销。

\subsection{算子融合}

\term{算子融合}{Operator Fusion}{}将多个相邻计算节点合并至同一个计算内核中执行，减少中间张量的显存读写与内核启动开销。以 $\vect{y}=\phi(\vect{x}\mat{W}+\vect{b})$ 为例，若矩阵乘法后将中间结果 $\mat{Z}=\vect{x}\mat{W}$ 完整写回显存，偏置加法与激活函数再分别发起全局内存读写，引入的显存访存量正比于 $\mat{Z}$ 的元素总数。将偏置相加与非线性激活作为矩阵乘法累加寄存器的尾部操作直接发射（Epilogue Fusion），中间张量完全驻留在片上寄存器中，消除了该部分显存带宽占用。

算子融合受硬件物理资源的严格约束：当中间特征需被多个下游分支复用时，融合可能导致重复计算；融合算子若引入过多局部变量，会导致寄存器溢出（Register Spill）至局部显存，显著增加访存延迟；跨线程块的全局归约或跨卡通信需要显式硬件屏障，限制了融合边界的扩展。因此融合设计需综合权衡片上数据复用深度、寄存器压力与硬件同步开销。

\subsection{图编译与执行回放}

\term{图编译}{Graph Compilation}{}将动态计算图静态化为针对目标设备指令级优化的执行程序。\term{执行图回放}{Execution Graph Replay}{}（如 CUDA Graphs）预先捕获并固化 GPU 任务的内核拓扑与内存依赖，后续推理只需单次触发即可由硬件工作引擎完成全图下发，消除了每次前向由主机 CPU 发起频繁内核发射引入的调度气泡。

若模型控制流中存在主机依据设备端张量值做动态分支判定的操作（如动态 Early-exit 或基于张量值的异常中断），主机与设备间的同步将打断图捕获链路。工程实现应严格分离数据依赖与元数据依赖，将循环展开与掩码长度控制收敛于静态张量维度内。

\subsection{形状特化的摊销模型}
设单次静态编译耗时为 $C$，动态解释执行单次耗时为 $t_e$，编译特化后单次执行耗时为 $t_c<t_e$。在调用频次为 $n$ 的推理服务中，图编译取得净时间收益的充要条件为
\begin{equation}
 C+nt_c<nt_e\quad\Longleftrightarrow\quad n>\frac{C}{t_e-t_c}.
 \label{eq:kernel-amortize}
\end{equation}
若 $C=\qty{3}{\second}$，单次执行耗时由 $t_e=\qty{2}{\milli\second}$ 降至 $t_c=\qty{1.5}{\milli\second}$，则至少需要 $n>\num{6000}$ 次调用才能摊平静态编译开销。在动态变长推理场景中，若对每个出现的新输入形状均发起即时特化编译，调用频次极低的冷门形状将无法回收编译成本，甚至引发显存膨胀与首字延迟劣化。

\term{形状分桶}{Shape Bucketing}{}将连续变化的输入长度离散化映射至一组预设的几何或算术边界桶。设输入序列长度随机变量为 $T$，映射至其所在的桶上界 $b(T)$，则平均填充无用词元数为 $\mathbb{E}[b(T)-T]$。对于稠密自注意力，填充引入的无效计算量正比于二阶矩 $\mathbb{E}[b(T)^2-T^2]$。分桶策略需要在编译缓存开销、预留显存碎片与填充算力损耗之间寻找权衡。

\section{推测解码}

\subsection{推测解码的分布保持}
\term{推测解码}{Speculative Decoding}{}由计算成本较低的草稿模型提出前向候选，再由参数量更大的目标模型执行并行因果验证\autocite{leviathan2023speculative}。本节重点讨论保证输出序列与目标自回归采样分布完全同一的精确验证算法。

\begin{assumption}[采样分布条件]
草稿分布 $q$ 与目标分布 $p$ 定义在同一离散词元空间 $\mathcal{V}$ 上，且满足全概率归一化。候选词元严格依条件分布 $q(\cdot\mid x_{<t})$ 独立采样，分布已包含温度、截断（Top-$k$/Top-$p$）等变换。草稿分布与目标分布基于相同的已确认历史前缀展开。
\end{assumption}

对于草稿提出的候选词元 $x\sim q$，当 $q(x)>0$ 时，定义其接受概率为
\begin{equation}
 a(x)=\min\left\{1,\frac{p(x)}{q(x)}\right\}.
\end{equation}
候选词元 $x$ 被提议且被目标模型接受的联合概率质量为 $q(x)a(x)=\min(p(x),q(x))$。全词表总接受率为
\begin{equation}
 A=\sum_{x\in\mathcal{V}}\min(p(x),q(x))
 =1-\frac{1}{2}\sum_{x\in\mathcal{V}}|p(x)-q(x)|=1-D_{\mathrm{TV}}(p,q).
 \label{eq:kernel-accept-tv}
\end{equation}
式\eqref{eq:kernel-accept-tv}表明，单步理论接受率等于 $1$ 减去目标分布与草稿分布的总变差距离（Total Variation Distance）。

若候选词元未通过接受检验，定义单步拒绝概率为 $Z=1-A$。此时算法不丢弃当前步，而是从校正残差分布
\begin{equation}
 r(x)=\frac{[p(x)-q(x)]_+}{Z}
 \label{eq:kernel-residual}
\end{equation}
中独立采样替代词元，其中 $[u]_+=\max(u,0)$。由 $\sum_x p(x)=\sum_x q(x)=1$ 可知正差累加和严格为 $\sum_x [p(x)-q(x)]_+=Z$，因此残差分布 $r(x)$ 满足合法概率分布的非负性与归一化条件。该机制下的实际输出概率为
\begin{equation}
 \Pr(\mathrm{output}=x)=q(x)a(x)+Z r(x)=\min(p(x),q(x))+[p(x)-q(x)]_+=p(x).
 \label{eq:kernel-exact-speculation}
\end{equation}
式\eqref{eq:kernel-exact-speculation}证明了输出分布对目标分布的严格精确保持。\keyconcept{推测解码通过接受概率 $a(x)=\min(1, \frac{p(x)}{q(x)})$ 与正残差重采样分布 $r(x)=\frac{[p(x)-q(x)]_+}{Z}$，在数学上严格精确保持目标自回归采样分布。}当 $Z=0$ 时，草稿分布与目标分布完全一致，候选必定被全额接受。对于草稿概率为零但目标概率非零的词元（即 $q(x)=0$ 且 $p(x)>0$），其在残差分布中的分子保留完整的 $p(x)$，在发生拒绝时可被按比例重采样恢复。因此推测解码不要求草稿分布必须具备全词表的全正支撑集。

\begin{example}
设词表大小为 $3$，目标概率分布为 $p=(0.5,0.3,0.2)$，草稿分布为 $q=(0.2,0.5,0.3)$。
三个词元的逐项接受概率为 $a=(\min(1,\frac{0.5}{0.2}),\min(1,\frac{0.3}{0.5}),\min(1,\frac{0.2}{0.3}))=(1,0.6,\frac{2}{3})$。
接受分支的联合概率质量为 $(0.2\times 1,0.5\times 0.6,0.3\times \frac{2}{3})=(0.2,0.3,0.2)$，总接受率 $A=0.2+0.3+0.2=0.7$。
拒绝概率 $Z=1-0.7=0.3$。正残差项 $[p-q]_+=(0.3,0,0)$，归一化残差分布为 $r=(1,0,0)$。
拒绝发生时必定采样输出第一个词元，提供修正质量 $0.3\times (1,0,0)=(0.3,0,0)$。
将接受分支与修正分支相加：$(0.2,0.3,0.2)+(0.3,0,0)=(0.5,0.3,0.2)=p$，精确还原目标分布。

若在发生拒绝后错误地直接重新从\emph{全目标分布} $p$ 采样，而非从\emph{残差分布} $r$ 采样，合并后的实际输出概率为 $(0.2,0.3,0.2)+0.3\times(0.5,0.3,0.2)=(0.35,0.39,0.26)\ne p$，破坏了原始自回归生成分布。
\end{example}

\subsection{候选验证的执行过程}

草稿模型在已确认前缀后自回归串行生成 $k$ 个候选词元 $x_1,\ldots,x_k$。目标模型将这 $k$ 个候选词元拼接至输入，利用因果掩码执行单次批量前向传播，同时获得这 $k$ 个词元各自上下文下的目标条件分布 $p(x_i\mid x_{<i})$ 以及全部候选接受后的第 $k+1$ 步分布。

验证过程必须沿时间步由前向后严格按序判别：仅当 $x_1,\ldots,x_{j-1}$ 全部通过接受测试时，$x_j$ 的验证环境才与目标模型的自回归因果前缀严格等价。若首个拒绝发生在位置 $j$，则保留已确认前缀 $x_{<j}$，在位置 $j$ 依残差分布 $r_j(x)$ 采样生成修正词元，并丢弃 $x_{>j}$。若 $k$ 个候选全部被接受，则直接从第 $k+1$ 步目标分布中额外采样一个词元。逐位置应用式\eqref{eq:kernel-exact-speculation}并通过数学归纳法，保证了推测解码生成的任意长度序列严格服从目标大模型的自回归联合分布。

\begin{figure}[htbp]
\centering
\begin{tikzpicture}[node distance=7mm,box/.style={draw,rounded corners,align=center,text width=31mm,minimum height=11mm},>=stealth]
\node[box] (prefix) {已确认前缀\\缓存检查点};
\node[box,right=of prefix] (draft) {草稿提出 $k$ 词元\\保存条件概率};
\node[box,right=of draft] (verify) {目标因果验证\\顺序接受判断};
\node[box,below=of verify] (commit) {提交接受前缀\\修正或额外词元};
\node[box,left=of commit] (rollback) {裁剪未接受缓存\\更新逻辑长度};
\draw[->] (prefix)--(draft);\draw[->] (draft)--(verify);\draw[->] (verify)--(commit);\draw[->] (commit)--(rollback);\draw[->] (rollback)--(prefix);
\end{tikzpicture}
\caption[推测解码提议与提交边界]{推测解码提议与提交边界。草稿模型生成的键值状态为暂存状态，验证拒绝后的后缀随即失效且不再参与注意力计算。}\label{fig:kernel-speculation}
\end{figure}

在显存管理层面，拒绝位置产生的修正词元并未经历目标模型的前向特征提取。因此系统须严格解耦“逻辑输出序列长度”与“显存物化 KV 长度”：未被目标模型前向计算的词元，其键值状态留待下一轮首步前向时统一物化。若仅在应用层回滚文本而未在底层物理缓存中裁剪被拒绝后缀的键值状态，残留的无效键值将在下一轮计算中污染因果注意力。

\begin{algorithm}[保持目标分布的推测解码]
\AlgoInput{目标模型 $p$、草稿模型 $q$、已确认前缀 $x_{\le t}$、每轮推测步数 $k$、终止条件。}
\AlgoOutput{符合目标模型分布的新增生成词元。}
\begin{enumerate}
\item 记录当前已确认前缀的 KV 缓存边界。草稿模型自回归生成 $k$ 个候选词元 $x_{1},\ldots,x_{k}$，并保存每一步的完整条件提议分布 $q_i(\cdot)$。
\item 目标模型将 $k$ 个候选一次性并行前向，获取各候选条件分布 $p_1,\ldots,p_k$ 及后续分布 $p_{k+1}$。
\item 按序遍历 $i=1,\ldots,k$：从均匀分布 $u\sim\mathcal{U}(0,1)$ 采样。若 $u\le a(x_i)$，接受并提交 $x_i$；若遇终止标记则提前退出。
\item 若在第 $j$ 步发生首次拒绝，按残差分布式\eqref{eq:kernel-residual}采样修正词元 $x'_j$ 并提交；丢弃后续未接受候选，回滚目标模型与草稿模型在位置 $j$ 之后的物理缓存。
\item 若 $k$ 个候选全部被接受且未触发终止，从 $p_{k+1}$ 采样一个新词元并提交。
\item 更新逻辑序列指针与物理缓存水位，进入下一轮迭代。
\end{enumerate}
\end{algorithm}

\subsection{推测解码的加速条件}

设一轮中每个候选词元在前面均被接受的前提下，条件接受概率为常数 $\alpha$。提交至少 $i+1$ 个有效词元要求前 $i$ 个词元均被接受，因此单轮平均提交词元数 $N$ 的数学期望为
\begin{equation}
 \mathbb E[N]=\sum_{i=0}^k\Pr(N\ge i+1)
 =\sum_{i=0}^k\alpha^i
 =\frac{1-\alpha^{k+1}}{1-\alpha},
 \label{eq:kernel-spec-speed}
\end{equation}
当 $\alpha=1$ 时式\eqref{eq:kernel-spec-speed}的极限值为 $k+1$。若各步接受概率随推测深度衰减，应代入各级联合生存概率计算；直接用平均接受率替代各阶幂次会高估期望提交量。

设草稿模型单步前向耗时为 $c_d$，目标模型单次验证 $k$ 个词元的并行前向耗时为 $c_v(k)$，缓存回滚与采样控制开销为 $c_m$，目标模型基准单步自回归耗时为 $c_t$。推测解码的实际端到端加速比为
\begin{equation}
 S\approx\frac{c_t\mathbb E[N]}{k c_d+c_v(k)+c_m}.
\end{equation}
以 $k=3,\alpha=0.8$ 为例，单轮期望产出词元数为 $\mathbb{E}[N]=2.952$。设 $c_t=\qty{10}{\milli\second}, c_d=\qty{1}{\milli\second}, c_v(3)=\qty{12}{\milli\second}, c_m=\qty{1}{\milli\second}$，端到端加速比为
\begin{equation*}
 S\approx\frac{10\times 2.952}{3\times 1+12+1}=\frac{29.52}{16}\approx 1.85.
\end{equation*}
若由于模型架构或批量过大导致验证耗时上升至 $c_v(3)=\qty{30}{\milli\second}$，则加速比降至 $\frac{29.52}{34}\approx 0.87<1$，推测解码反而劣化了系统性能。因此推测加速成立的前提是：草稿前向与目标验证的合计开销显著低于目标模型串行解码相同数量词元的开销。

表\ref{tab:speculative-decoding}比较了几种推测解码方法的草稿来源、额外存储和训练需求。选择方法时还需在目标任务上测量接受率及验证成本。

\begin{table}[htbp]
\centering
\caption{推测解码主要技术范式对比}\label{tab:speculative-decoding}
\begin{tabularx}{\linewidth}{@{}lXXXX@{}}
\toprule
范式分类 & 代表方案 & 草稿生成机制 & 额外存储与训练开销 & 适用场景与限制 \\
\midrule
独立小模型推测 & Speculative Decoding\autocite{leviathan2023speculative} & 同词表的轻量预训练小模型自回归生成 & 需额外加载小模型权重与独立 KV 缓存 & 通用场景，但跨模型词表/对齐偏差制约接受率 \\
多头自投机 & Medusa, Blockwise & 主模型顶层挂接多个前向预测头，树状注意力验证 & 仅需微调少量轻量头参数，无额外主干显存 & 保持单模型部署；单步独立预测忽略词元间依赖 \\
特征级推测 & EAGLE, EAGLE-2 & 结合顶层隐藏特征与词元嵌入，单层 Transformer 前向 & 需训练单个解码层，显存占用极低 & 结合上下文语义，接受率显著优于独立小模型 \\
非参数检索匹配 & Prompt Lookup, REST & 基于当前输入上下文或离线检索库进行 $n$-gram 匹配 & 无需额外参数与模型，零显存开销 & 代码编辑、多轮文档问答等高复现文本场景 \\
\bottomrule
\end{tabularx}
\end{table}

\term{滑动窗口注意力}{Sliding-Window Attention}{}通过将键值可见范围截断至最近的 $W$ 个词元来降低计算与显存占用。对于未在训练阶段引入窗口因果掩码的全注意力模型，运行时截断窗口之外的注意力质量将导致输出表征偏离真实值。设被截断位置的原注意力权重和为 $\delta$，所有值向量的二范数满足 $\|v_j\|_2\le M$。令保留区域重新归一化后的输出为 $\vect{o}'$，在 $\delta<1$ 的条件下满足误差上界：
\begin{equation}
 \|\vect{o}-\vect{o}'\|_2\le 2\delta M.
\end{equation}
该上界由将原输出分解为 $(1-\delta)\vect{o}'+\delta\vect{o}_{\mathrm{drop}}$ 并结合三角不等式推导得出。表征偏离幅度完全取决于被丢弃位置的注意力概率质量 $\delta$；仅当远距离历史的注意力权重衰减至极小值时，截断引入的表征误差才受控。

推理优化的系统验证需在不同抽象层级设立正交的标准：算子级核验单元测试中的绝对/相对数值误差容限（ATol/RTol）；状态级核验增量解码与全量前向的键值缓存一致性；算法级核验推测解码输出概率分布的经验频率与目标模型理论分布的拟合度；服务级核验首次词元延迟（TTFT）、词元生成吞吐（TPS）以及并发排队下的长尾时延（P99）。区分冷启动特化编译与稳态执行流水线，避免以微基准内核峰值替代端到端工程表现。




\chapterreferences
\myendchapterdocument
